\documentclass{article}

\usepackage{booktabs}
\usepackage{multirow}
\usepackage{amsmath}
\usepackage{hyperref}
\usepackage{overpic}
\usepackage{amssymb}

\usepackage[accepted]{dlaiml2026}

\dlaititlerunning{Document-Aware KV Caching for Long-Document QA}

\begin{document}

\twocolumn[
\dlaititle{Document-Aware KV-Cache Management for Repeated Long-Document Question Answering}

\begin{center}September 2026\end{center}

\begin{dlaiauthorlist}
\dlaiauthor{[Student name]}{}
\end{dlaiauthorlist}

\dlaicorrespondingauthor{[Student name]}{[institutional email]}

\vskip 0.3in
]

\printAffiliationsAndNotice{}

\begin{abstract}
Answering several questions about one document makes a language model process
much of the same input repeatedly. I study whether caching each document's
KV state as one entry can reduce this work and simplify cache management.
I compare document, fixed-block and radix caches with Qwen2.5-1.5B on
QuALITY, where each question supplies its article. Over 2,086 requests, document/LRU
reduces mean time to first token (TTFT) by 21.9\% with random ordering and
63.5\% with Zipf traffic. CPU INT8 increases capacity but changes 12 answers
on the random trace. In smaller experiments,
a page arena does not improve latency in the Transformers runner, while a
Triton kernel speeds up INT8 restoration by $1.209\times$ after warm-up.
\end{abstract}

\section{Introduction and scope}

An assistant answering questions about a long article receives nearly the
same prompt each time: the article stays, while the question changes, and
Transformers retain attention keys and values during generation
\cite{vaswani2017attention}. Reusing this state across requests with an exact
common prefix avoids processing the article again, but keeping every article's
KV state may not fit in memory, so some entries must be evicted.

I initially proposed studying this in RAG \cite{lewis2020rag}: requests would
use two or three documents, with the same combinations recurring across
questions. A company legal assistant might repeatedly need the same statutory
provisions and internal policy. I did not find a ready-made public serving
trace with this pattern, and constructing one would require choosing document
combinations, their order and repetition frequency. I therefore use QuALITY,
where each question supplies its \texttt{article\_id} \cite{pang2022quality},
to study reuse of one known document first. Retrieval and multi-document
prompts are outside this experiment.

\section{Method}

Every prompt uses the same order:
\[
 \underbrace{\text{system prefix}}_{L0}\;\Vert\;
 \underbrace{\text{article}}_{\text{reusable}}\;\Vert\;
 \underbrace{\text{question + A/B/C/D}}_{\text{request-specific}}.
\]
L0 is never evicted or counted as article reuse; questions and options are
never cached. For $T$ tokens, $L$ layers,
$H_{kv}$ KV heads, head width $D$ and $b$ bytes per value, the KV size is
$B_{KV}(T)=2LH_{kv}DTb$. Qwen2.5-1.5B has 28 layers, two KV heads and width
128: each FP16 token therefore costs 28 KiB \cite{yang2024qwen25}.

All three caches use the same Transformers runner. \textbf{Document} stores and
evicts whole articles, \textbf{fixed-block} uses 256-token blocks inspired by
vLLM \cite{kwon2023pagedattention}, and \textbf{radix} uses longest-prefix lookup
inspired by SGLang \cite{zheng2024sglang}. These are local implementations, not
benchmarks of either engine. The reported runs use LRU, with GDSF added for Zipf.
KV tensors use GPU FP16 or symmetric CPU INT8 with per-layer, per-KV-head
scales. Restoring CPU INT8 entries requires transfer and dequantization, so
TTFT includes both costs.

To compare how much useful work each cache avoids, I count article tokens:
\begin{equation}
 h_{article}=\frac{\sum_i\text{restored article tokens}_i}
 {\sum_i\text{requested article tokens}_i}.
\end{equation}
The no-cache baseline processes the same system, article and suffix stages,
but discards article KV after every request. Using this segmented baseline
matters: otherwise, differences between one full forward pass and several
smaller passes could be mistaken for a benefit of caching.

\section{Experimental protocol}

The data is HTML-stripped QuALITY v1.0.1. Merging the two writer records per
article preserves all 381 articles and 6,737 questions. Accuracy is measured
on dev; the test split is used only for cache traces because its labels are
withheld. With seed 42, random traffic shuffles the questions, while Zipf
favors popular articles and cycles through their questions. Answers are
selected by scoring A/B/C/D, with sequence-likelihood
fallback when a label is not a single token.

I ran twelve Qwen2.5-1.5B-Instruct configurations on a Tesla T4. Each serves
2,086 aligned requests, with 4 GiB for cached runs. Random visits every dev
question once; Zipf contains 911 distinct Q\&A, so I also remove duplicates
when evaluating accuracy. A 108-run no-inference matrix selected the cache
strategies and policies. Code and experiment details are at
\url{https://github.com/i0nut02/rag-kvcache}.

\section{Results}

\begin{table}[t]
\caption{Selected full-dev inference paths: 2,086 requests per path, 4 GiB for cached paths.
TTFT mean and p90 are in seconds; speedup compares means with the aligned
segmented control. $h$ excludes L0. Storage is accelerator FP16 except CPU
INT8/Triton.}
\label{tab:main}
\centering
\begin{footnotesize}
\setlength{\tabcolsep}{3pt}
\begin{tabular}{llrrrr}
\toprule
Traffic & Cache & $h_{article}$ & Mean & p90 & Speedup \\
\midrule
Random & Segmented & 0.00\% & 1.843 & 2.623 & $1.000\times$ \\
Random & Document/LRU & 22.94\% & 1.439 & 2.562 & $1.281\times$ \\
Random & Fixed-256/LRU & 23.15\% & 1.474 & 2.605 & $1.251\times$ \\
Random & Radix/LRU & 22.77\% & 1.528 & 2.736 & $1.206\times$ \\
Random & Document/INT8 & 44.67\% & 1.168 & 2.752 & $1.578\times$ \\
\midrule
Zipf & Segmented & 0.00\% & 2.208 & 2.670 & $1.000\times$ \\
Zipf & Document/LRU & 65.38\% & 0.806 & 2.543 & $2.741\times$ \\
Zipf & Fixed-256/LRU & 65.27\% & 0.863 & 2.566 & $2.558\times$ \\
Zipf & Radix/LRU & 65.23\% & 0.762 & 2.457 & $2.898\times$ \\
Zipf & Document/GDSF & 71.46\% & 0.633 & 2.310 & $3.487\times$ \\
Zipf & Document/INT8 & 81.66\% & 0.455 & 2.207 & $4.851\times$ \\
\bottomrule
\end{tabular}
\end{footnotesize}
\end{table}

\paragraph*{Cache organization.}
Table~\ref{tab:main} shows that caching reduces average latency in both
workloads, with larger savings under Zipf because popular articles are
requested repeatedly. The choice of cache structure makes a smaller
difference: document and radix differ by less than 0.5\% in mean TTFT across
the repeated runs, whereas fixed-256 is slower in every repetition.
Document's clearer advantage is the amount of bookkeeping it avoids by
managing each article as one entry, using about 22 KB of metadata compared
with 5.72 MB for fixed-256 on full random dev.

These comparisons also show why a partial hit is useful only in proportion
to how much text it restores. On Zipf, the median partial hit covers 6,400
article tokens with fixed blocks but only one token with radix, so counting
both as a hit would hide a large difference in avoided prefill work.
The token-weighted metric captures this distinction, while p90 shows the
remaining cost of misses, which still require processing most of the article.
Eviction policy matters when some articles are more popular than others:
on Zipf, GDSF raises article-token reuse from 65.38\% to 71.46\% and lowers
mean TTFT by 21.4\% compared with document/LRU.

\paragraph*{INT8 and answer accuracy.}
INT8 trades numerical precision for enough space to keep more articles,
with 4 GiB nearly matching the capacity of 8 GiB of FP16 in simulation.
On random traffic this improves mean latency, though p90 becomes worse,
and unchanged overall accuracy hides 12 answer changes: six become correct
and six become wrong. Zipf needs particular care because one helpful change
is counted 32 times as the same question recurs; counting each unique Q\&A
once instead gives 496 correct answers for INT8 versus 499 for the baseline,
out of 911. Across these unique items, 14 labels change: seven correct-to-wrong,
four wrong-to-correct and three wrong-to-different-wrong. The useful gain is
therefore cache capacity, not better answer quality, even when aggregate
accuracy appears unchanged or higher.

\paragraph*{Arena and Triton experiments (100 requests).}
The arena and Triton experiments illustrate how an optimization of cache
storage does not necessarily improve the whole runner. In random-dev tests
with the same model, T4 and seed, both arena page sizes increase latency
because Transformers still requires contiguous KV tensors to be rebuilt
before use. Triton has a more direct benefit in INT8 restoration
\cite{tillet2019triton}: at 31 matched cache-hit positions, complete restore
falls from 37.45 to 30.98 ms with identical scores to PyTorch INT8.
That saving applies after warm-up and must first recover the 2.671 s JIT
startup cost. At the observed 31\% hit rate, the estimated break-even point
is roughly 1,300 requests, using a simple extrapolation detailed in the
repository rather than a measured crossover.

Full run-level tables, three-repetition ranges, accuracy breakdowns and the
arena/Triton microbenchmarks are available in \path{docs/results.md} and
\path{docs/generated/} at the repository URL above.

\section{Discussion and conclusion}

For this workload, document/LRU combines low management cost with latency
close to radix, while fixed-block is consistently slower after tuning. The
evidence covers QuALITY, two Qwen sizes and one T4; full configurations ran
once and shorter comparisons ran three times on one seed, so they show timing
variation rather than a general ranking. RAG would also require versioned
retrieved prompts and separate retrieval timing.

\clearpage

\section{Statement on the use of AI}

The main ideas and initial project structure were my own. I used generative
AI extensively to develop the implementation: the Triton kernel, experiment
scripts and much of the supporting test and analysis code were mainly
AI-written. AI also helped me refine the design, read the literature and
refactor code.

I used AI to prepare tables and plots from saved results, discuss their
interpretation, check cache correctness and revise this report. This included
spotting unfair timing comparisons and planning further tests.
I reviewed the generated code and checked the tests and saved outputs
during development.

\bibliography{references.bib}
\bibliographystyle{dlaiml2026}
\end{document}
